\section{Introduction}
Two plausible diagnoses can each prescribe a consistent repair while their token-level combination prescribes neither. Consider an illustrative two-token action space in which one component assigns probability one to \texttt{AA} and another assigns probability one to \texttt{BB}. An equally weighted mixture of their complete sequences assigns probability one half to each coherent action and zero to \texttt{AB} or \texttt{BA}. If the component is independently remixed at each token, all four strings receive probability one quarter. This arithmetic example is not an experimental result: it isolates the distinction between preserving a hypothesis through a continuation and repeatedly forgetting which hypothesis was selected.

Program repair makes this distinction relevant because sparse execution evidence can support incompatible changes to control flow, indexing, or input interpretation. Generating more candidates explores possible answers, but does not by itself identify what evidence supports each explanation. Choosing only the most probable hypothesis can discard an alternative that is still plausible; averaging latent hypotheses need not produce a meaningful diagnosis. The appropriate action may therefore require both a calibrated multimodal belief and a mechanism that preserves one component long enough to express a complete edit.

We propose Particle-Belief Program Filtering (PBPF), a candidate-specific belief model paired with continuation-level conditioning. For unchanged source code, ordered execution outcomes reweight fixed failure particles. A patch creates a child latent sampled from an evidence-conditioned proposal, with an explicit importance correction relative to the learned transition. Parent and child accounting is kept separate from unrelated candidate slots. This connects sequential Monte Carlo and learned filtering objectives~\citep{delmoral2006,maddison2017,naesseth2018} to a diagnosis variable rather than to the generated token trajectory.

To act, PBPF samples one posterior component and holds its projected soft prefix fixed for the entire repair. The actor objective is a mixture of complete sequence likelihoods. Our comparison with rollout-level particle methods~\citep{roulette} concerns what a particle denotes, not whether probabilistic inference is new. Likewise, feedback-conditioned repair~\citep{selfdebug} and candidate allocation~\citep{rex} are complementary controls: they test whether the explicit belief contributes beyond transcript access and search policy.

Preserving multimodality is only useful if the modes predict something beyond the observations already seen. The prospective protocol therefore first evaluates evaluator-withheld outcomes, then tests repair under identical initial candidate hashes and generation allowances. All outcomes are pending. Three proposed contributions map to these distinct tests:
\begin{itemize}
\item A sequential latent diagnosis formulation whose piecewise-static semantics and importance weights can be checked against exact finite posteriors before code-repair claims are considered.
\item A whole-sequence posterior mixture with sample-once generation, together with explicit MAP, mean, token-remix, and full-ensemble comparisons that isolate coherence from extra decoding.
\item A gated predictive-and-repair protocol using calibrated future-outcome scoring, information-matched baselines, evaluator isolation, and source-cluster confidence intervals to test whether uncertainty improves action.
\end{itemize}
